\documentclass[10pt,a4paper]{amsart}
\usepackage[margin=19mm]{geometry}
\usepackage{amsmath,amssymb,mathtools}
\usepackage[scr=rsfs,cal=euler]{mathalfa}
\usepackage{xfrac}
\usepackage[unicode,hidelinks,bookmarks=false]{hyperref}
\newcommand{\ie}{i.e.\ }
\newcommand{\cf}{cf.\ }
\newcommand{\resp}{respectively\ }
\newcommand{\Subsection}{\subsection}
\providecommand{\nobreakdash}{\nobreak\hbox{-}}
\setlength{\emergencystretch}{2em}
\multlinegap0pt
\usepackage{bm}
\usepackage{bbm}
\usepackage{dsfont}
\usepackage{xspace}
\usepackage{cancel}
\usepackage{mathtools}
\usepackage{amsthm}
\makeatletter
\@ifundefined{equation}{\newtheorem{equation}{Equation}}{}
\@ifundefined{prop}{\newtheorem{prop}[equation]{Proposition}}{}
\makeatother
\makeatletter
\expandafter\ifx\csname AMS\endcsname\relax
\newenvironment{AMS}{}{}
\fi
\newcommand{\Bm}[1]{ \boldsymbol{ #1 }}
\newcommand{\C}{\mathbb{C}}
\newcommand{\N}{\mathbb{N}}
\newcommand{\Om}{\Omega}
\newcommand{\bei}{\begin{itemize}}
\newcommand{\beqan}{\[\begin{aligned}}
\newcommand{\bk}{\backslash}
\newcommand{\cc}{\circ}
\newcommand{\cd}{\cdot}
\newcommand{\de}{\delta}
\newcommand{\eei}{\end{itemize}}
\newcommand{\eeqan}{\end{aligned}\]}
\newcommand{\ep}{\varepsilon}
\newcommand{\ga}{\gamma}
\expandafter\ifx\csname keywords\endcsname\relax
\newenvironment{keywords}{}{}
\fi
\newcommand{\la}{\lambda}
\newcommand{\lb}{\left\lbrace}
\newcommand{\lo}{\longrightarrow}
\newcommand{\mI}{\mathcal{I}}
\newcommand{\mM}{\mathcal{M}}
\newcommand{\mO}{\mathcal{O}}
\newcommand{\mT}{\mathcal{T}}
\newcommand{\mZ}{\mathcal{Z}}
\newcommand{\mt}{\mapsto}
\newcommand{\ov}[1]{ \overline{#1}}
\newcommand{\q}{\quad}
\newcommand{\ra}{\rightarrow}
\newcommand{\rb}{\right\rbrace }
\newcommand{\sL}{\mathscr{L}}
\newcommand{\sig}{\sigma}
\newcommand{\sse}{\subseteq}
\newcommand{\ty}{\infty} 
\makeatother
\title[Multiplication Operator Semigroups on Banach lattice valued ]{Multiplication Operator Semigroups on Banach lattice valued continuous function spaces}
\author{Tobi David Olabiyi}
\date{Source: arXiv:2510.03607v1 (CC BY 4.0)}
\begin{document}
\maketitle

\noindent\emph{Source and licence.} Multiplication Operator Semigroups on Banach lattice valued continuous function spaces. Version arXiv:2510.03607v1 (CC BY 4.0), distributed under the Creative Commons Attribution 4.0 International licence, as recorded in the arXiv licence metadata for this exact version and shown on the abstract page https://arxiv.org/abs/2510.03607v1. Full rights notice: licenses/arxiv-2510.03607-CC-BY-4.0.txt.\par\smallskip

\noindent\textit{Editorial scope.} Counted unit: the Prop whose source label is \texttt{prop:mop1} in the article (source0.tex lines 261--277, source label \texttt{prop:mop1}), delivered together with the article's own complete original proof of it (source0.tex lines 280--329, one \texttt{proof} environment of 50 lines). No mathematics is written, repaired or re-derived: statement and proof are the article's own text line for line. Required context retained uncounted: none. Every definition, display and label of those lines is retained unchanged. Omitted from this package and not redistributed: 1--260, 278--279, 330--531. Nothing inside a retained excerpt is omitted or altered: every retained body is delivered exactly as the article's lines are written, so no omission receipt of any kind accompanies this package. Citations: the retained excerpts carry no citation command at all, so no bibliography entry of the article is transcribed and no reference is redistributed. Reference adaptation: none was needed and none was made; every reference inside the delivered unit resolves inside it, so no unresolved reference or citation is produced. Printed slips kept literal: no printed word, symbol, hypothesis or number of the counted statement or of its proof was altered. Layout and class: the article's own class is \texttt{article} with packages including graphicx, amsmath, amsfonts, amssymb, tikz, tikz-cd, amsthm, mathrsfs, xcolor, hyperref; the wrapper is the lane's pinned \texttt{amsart} editorial wrapper at 10pt on A4, which restates the theorem-like environments the retained text needs and adds the article's own macro definitions that the retained text needs (\texttt{\textbackslash Bm}, \texttt{\textbackslash C}, \texttt{\textbackslash N}, \texttt{\textbackslash Om}, \texttt{\textbackslash bei}, \texttt{\textbackslash beqan}, \texttt{\textbackslash bk}, \texttt{\textbackslash cc}, \texttt{\textbackslash cd}, \texttt{\textbackslash de}, \texttt{\textbackslash eei}, \texttt{\textbackslash eeqan}, \texttt{\textbackslash ep}, \texttt{\textbackslash ga}, \texttt{\textbackslash la}, \texttt{\textbackslash lb}, \texttt{\textbackslash lo}, \texttt{\textbackslash mI}, \texttt{\textbackslash mM}, \texttt{\textbackslash mO}, \texttt{\textbackslash mT}, \texttt{\textbackslash mZ}, \texttt{\textbackslash mt}, \texttt{\textbackslash ov}, \texttt{\textbackslash q}, \texttt{\textbackslash ra}, \texttt{\textbackslash rb}, \texttt{\textbackslash sL}, \texttt{\textbackslash sig}, \texttt{\textbackslash sse}, \texttt{\textbackslash ty}), with extra wrapper packages (bm, bbm, dsfont, xspace, cancel, mathtools). The delivered numbering is the unit's own. Assets: no retained span contains a figure environment, a graphics-inclusion command, a PDF-page inclusion, a table or a TikZ picture; no image file is copied into this package, the package declares no asset dependency and no image file is redistributed with it. Licence: version arXiv:2510.03607v1 of this article is distributed under the Creative Commons Attribution 4.0 International licence, as recorded in the arXiv licence metadata for this exact version and shown on the abstract page https://arxiv.org/abs/2510.03607v1. A full rights notice is delivered at licenses/arxiv-2510.03607-CC-BY-4.0.txt. Characters: every retained excerpt is pure ASCII, so the delivered engine, XeLaTeX with its default Latin Modern text font, reports no missing character.
\begin{prop}\label{prop:mop1}
	Let $\mM_\phi : D(\mM_\phi)\sse C_0(\Om, E) \lo C_0(\Om, E); s \mt \phi\cc s$ be the multiplication operator induced by a continuous function $\phi : \Om \lo  \Bm{\mZ}(E)_\textbf{s}$. Then, we have the following.
	\bei	
	\item[(i)] The operator $(\mM_\phi,  D(\mM_\phi)  )$ is closed and densely defined.
	
	\item[(ii)] The operator $\mM_\phi$ is bounded (with $D(\mM_\phi) =C_0(\Om, E)$) if and only if $\phi$ is bounded, i.e., $\phi\in C_b(\Om, \Bm{\mZ}(E)_\textbf{s})$. In this case,  $\mM_\phi\in\Bm{\mZ} ( C_0(\Om, E))$ and one has
	$$||\mM_\phi||= ||\phi||= \sup_{x\in\Om}||\phi(x)||_{\sL(E)}.$$
	
	\item[(iii)] The operator $\mM_\phi$ has a bounded inverse if and only if $\phi$ has a bounded inverse $\phi^{-1}\in C_b(\Om, \Bm{\mZ}(E)_\textbf{s})$, i.e., $\phi(x)$ is invertible in $\Bm{\mZ}(E)_\textbf{s}$ for all $x\in\Om$, and the function $ x  \overset{\phi^{-1}}{\longmapsto}       \phi(x)^{-1}$ is in $C_b(\Om, \Bm{\mZ}(E)_\textbf{s})$.
	 In this case, one has
	\[  \mM_\phi^{-1}= \mM_{\phi^{-1} }.\]
	\item[(iv)] For the spectrum of $\mM_\phi$, we have that
		\[\sig(\mM_\phi)=  \Delta_\phi:=  \lb  \la\in\C :\la \in\bigcup_{x\in \Om}\sig(\phi(x)) \ \text{ or } \     \sup_{x\in\Om}|| R(\la, \phi(x))||_{\sL(E)}=\ty  \rb.\] 
	That is, the resolvent set of $\mM_\phi$ is given by 
	\[ \rho(\mM_\phi) = \Delta_u:= \lb \la\in\C : \la \in \bigcap_{x\in \Om} \rho(\phi(x)) \ \text{ and } \  \sup_{x\in\Om}|| R(\la, \phi(x))||_{\sL(E)}<\ty \rb.\]	
	\eei
\end{prop}

\begin{proof}
	\bei
	\item[(i)] It is clear that, the domain $D(\mM_\phi) $ always contain the space
	\[   C_c(\Om, E):=  \lb s\in C(\Om, E) : \text{supp }s \text{ is compact} \rb\] of all continuous functions $s : \Om \lo E$ having compact support
	\[\text{supp }s:=  \ov{ \lb x\in\Om : s(x)\ne 0\in E \rb}.  \]
	So, to prove that the operator $(\mM_\phi,  D(\mM_\phi)  )$  is densely defined, it suffices to show $C_c(\Om, E)$ is norm dense in $C_0(\Om, E).$ Indeed, if we take approximate identity $(e_j)_{j}$ of $C_0(\Om)$, and without loss of generality, we may assume that $e_j$ has compact support for each $j$, so that for every $s\in C_0(\Om, E)$, we have that $e_j\cd s\in C_c(\Om, E)$ for each $i$, and $\lim_{j}||s - e_j\cd s|| = 0$.
	
	
	
	Next, we show that $(\mM_\phi,  D(\mM_\phi)  )$ is a closed operator. Let $(s_n)_{n\in\N}\sse D(\mM_\phi)$ be a sequence converging to $s\in 	C_0(\Om, E)$ such that $\lim_{n\ra \ty} \phi \cc s_n =:s_0\in C_0(\Om, E)$ exists. Then, it must be that, for each $x\in \Om$, $\lim_{n\ra \ty}||s_n(x)-s(x)||_E=0$ and $\lim_{n\ra \ty}||\phi(x)s_n(x) - s_0(x)||_E=0$, from which it follows that $\lim_{n\ra \ty}||\phi(x)s_n(x)-\phi(x)s(x)||_E=0$. Thus, for each $x\in \Om$, $s_0(x)=\phi(x)s(x)$, so that $s\in D(\mM_\phi)$ and $s_0= \phi\cc s$.
	\item[(ii)] If $\phi\in C_b(\Om, \Bm{\mZ}(E)_\textbf{s})$, we have that
	\[  ||\mM_\phi s|| = \sup_{x\in\Om}||\phi(x)s(x)||_E \le ||\phi||||s||   \]       
	for every $s\in C_0(\Om, E)$; hence $\mM_\phi$ is bounded with $||\mM_\phi||\le ||\phi||$. And on the other hand, assume $\mM_\phi$ is bounded. Now, for each $x\in\Om$, let $f_x$ be a continuous function in $C_0(\Om)$ with compact support satisfying $||f_x||= 1 =f_x(x).$ And for (fixed) $z_0\in E$ with $||z_o||=1$, consider, for each $x\in\Om$, the continuous mapping $f_x\otimes z_0 : \Om \lo E ; y \mt (f_x\otimes z_0)(y):= f_x(y)z_0$. Then, we have that  $f_x\otimes z_0 \in C_c(\Om, E)$ with $||f_x\otimes z_0||=1$; and moreover,  for each $x\in\Om$, 
	\[||\phi(x)z_0||= ||\mM_\phi(f_x\otimes z_0  )(x)||\le ||\mM_\phi(f_x\otimes z_0  )||\le ||\mM_\phi||\]
	implies that $||\phi(x)z||\le ||\mM_\phi||$ for every $z\in E$ with $||z||=1$. And as such, we have that
	$||\phi(x)||\le ||\mM_\phi||$ for every $x\in \Om$; hence, $\phi\in C_b(\Om, \Bm{\mZ}(E)_\textbf{s})$ with $$||\phi||=\sup_{x\in\Om}||\phi(x)||_{\sL(E)}\le ||\mM_\phi||.$$
	The fact that $\mM_\phi$ is a multiplication operator on 	$C_0(\Om, E)$ follows, since the pointwise $(\phi(x)_{x\in\Om})$ are multiplication operators on $E$ which are continuously bounded on $\Om$. Indeed, for each $s\in C_0(\Om, E)$,
		\beqan
	|\mM_\phi s|(x)	 =    &  |\phi(x)s(x)| \\
	\le & ||\phi(x)||_{\sL(E)}|s(x)|\\
	\le & ||\phi|| \cd |s|(x)
	\eeqan
	for all $x\in\Om$, implies that $|\mM_\phi s|\le ||\phi||\cd |s|$. Thus $\mM_\phi\in\Bm{\mZ} ( C_0(\Om, E))$.
	
	

	\item[(iii)] Suppose $\phi$ has bounded inverse $\phi^{-1}\in C_b(\Om, \Bm{\mZ}(E)_\textbf{s})$. This implies that $\phi^{-1}\cc\phi =\phi\cc\phi^{-1} =e \in C_b(\Om, \Bm{\mZ}(E)_\textbf{s})$.  Then, the induced (bounded) multiplication operator $$\mM_{\phi^{-1}} : C_0(\Om, E) \lo C_0(\Om, E) ; s \mt \phi^{-1}\cc s$$ is the inverse of $\mM_\phi$, since  $\mM_{\phi^{-1}}\mM_\phi s = s$ for all $s\in D(\mM_\phi)$ and	
	$   \mM_\phi \mM_{\phi^{-1}} = \mM_e = \mI\in \Bm{\mZ} ( C_0(\Om, E))$. \bigskip
	
		
	Now, suppose $\mM_\phi$ has bounded inverse, say $\mT$, i.e., $\mT\mM_\phi s = s$ for all $s\in D(\mM_\phi)$ and $\mM_\phi\mT = \mM_e = \mI$. Then, we obtain
\[||s|| \le ||\mT|| \cd ||\mM_\phi s|| \q \text{ for all } s\in D(\mM_\phi),\]
	whence, the estimate
\[\de:= \frac{1}{||\mT||}\le \sup_{x\in\Om}||\phi(x)s(x)||_E \q \text{ for all } s\in D(\mM_\phi), \ ||s||=1.\]
	
	And suppose on the contrary, without loss of generality, we assume either $\inf_{x\in\Om}\inf_{||z||=1}||\phi(x)z||	<\frac{\de}{2}$ or $\inf_{x\in\Om} \frac{1}{||\phi(x)^{-1}||}<\frac{\de}{2}$; then we can find an open set $\mO \sse \Om$ such that either $\inf_{||z||=1}||\phi(x)z||<\frac{\de}{2}$ or $\frac{2}{\de}<\sup_{||z||=1}||\phi(x)^{-1}z||$ for all $x\in\mO$ respectively. Let $f_0$ be a continuous function in $C_0(\Om)$ such that $||f_0||=1$ and  $f_0(x)=0$ for all $x\in\Om\bk\mO$, and consider, for any $z\in E$ with $||z||=1$, the continuous mappings $f_0\otimes z : \Om \lo E ; x \mt (f_0\otimes z)(x):= f_0(x)z$; so that $f_0\otimes z \in D(\mM_\phi)$ and $||f_0\otimes z||=1$. Then, the assumptions imply either $\sup_{x\in\Om} ||\phi(x)(f_0\otimes z)(x)||\le\frac{\de}{2}$ or $\frac{2}{\de} >\sup_{x\in\Om}||\phi(x)^{-1}(f_0\otimes z)(x) ||   =\sup_{x\in\Om}||\mT(f_0\otimes z)(x)||  $ for some $z\in E$ respectively, which is contradicting the above estimate. 
	
	
	Hence, we can conclude that each $\phi(x)$ is invertible for all $x\in\Om$ and $\sup_{x\in\Om}||\phi(x)^{-1}||= ||\mT || $, which would implies $\phi$ has bounded inverse $\phi^{-1}\in C_b(\Om, \Bm{\mZ}(E)_\textbf{s})$; so that $\mM_{\phi^{-1}}$ is the bounded multiplication operator on  $C_0(\Om, E)$ which coincides with the bounded inverse $\mT$.	
	
	Equivalently, $\Bm{\mZ} ( C_0(\Om, E))$ is an inverse-closed subalgebra of  $\sL( C_0(\Om, E)) $ and $ \mT(\phi \cc s) = s = \phi \cc \mT s$ for all $s\in D(\mM_\phi)$; it must be that the bounded inverse $\mT$ is a bounded multiplication operator on  $C_0(\Om, E)$, i.e., there exists $\psi\in C_b(\Om, \Bm{\mZ}(E)_\textbf{s})$ such that $\mM_\psi = \mT$. 	It, thus, follows that $\phi \cc \psi = e=\psi \cc\phi $, so that $\phi$ has bounded inverse $\phi^{-1}:=\psi$.
	
	\item[(iv)] By definition, one has $\la\in \sig(\mM_\phi)$ if and only if $\la - \mM_\phi = \mM_{\la -\phi}$ does not have bounded inverse. Thus, by (iii) above, it suffices to show that $\la -\phi$ has bounded inverse if and only if $\la\in\Delta_u$. The forward implication is clear, that  is if $\la -\phi$ has inverse  $ x  \overset{{(\la-\phi)}^{-1}}{\longmapsto} {(\la -\phi(x))}^{-1}$ in $C_b(\Om, \Bm{\mZ}(E)_\textbf{s})$, then $\la\in\Delta_u$. For the reverse, suppose $\la\in\Delta_u$ and consider the mapping $R: \Om \lo \Bm{\mZ}(E)_\textbf{s}; x \mt R(x):=R(\la,\phi(x) )= (\la -\phi(x))^{-1}$. Then, we see that $R$ is the pointwise inverse of $\la -\phi$  and is bounded. To prove that $R$ is continuous, let $x_\ga$ be a net in $\Om$ converging to $x$ in $\Om$. For all $\ep>0$ and each $z\in E$, using the continuity of $\la -\phi$, we eventually have $||z - (\la -\phi(x_\ga))R(x)z||<\ep$. As such, for each $z\in E$, we obtain
		\beqan
	||R(x_\ga)z - R(x)z||	 =    &   || R(x_\ga)[z - (\la -\phi(x_\ga))R(x)z] || \\
	\le & ||R(x_\ga)||_{\sL(E)}\ep
	\eeqan
	eventually, which can be made small as much as desired, since $R$ is bounded, i.e., $||R||:= \sup_{x\in\Om}||R(x)||_{\sL(E)}<\ty$. Thus, $R\in C_b(\Om, \Bm{\mZ}(E)_\textbf{s})$ which is the desired bounded inverse of $\la -\phi$.
	\eei
\end{proof}

\end{document}
